\section{Conclusion}\label{sec:conclusion}

This paper estimates an asymmetric IPV model of Brazilian electronic procurement and uses it to decompose the price effect of an exogenous 2018 extension of the SME-only rule into intensive and entry margins, and to price the welfare cost once the marginal cost of public funds is made explicit. Three findings stand out.

First, the intensive margin is the dominant channel: 70--80\% of the 23--33\% price increase comes from each bidder bidding less aggressively under the restricted pool, not from who shows up. The entry response runs the other way---SMEs roughly double their participation rate---and provides 20--30\% of offsetting relief. A Turnbull NPMLE ratifies the point estimate within 5\% in both classes.

Second, the welfare cost of the full set-aside is 28\% of procurement value in non-pharma (95\% CI 21--35\%) and 46\% in pharma (95\% CI 35--56\%) at an MCPF of 0.30. Two-thirds of that loss is allocative; one-third is tax distortion. The estimate is two to four times the 12\% benchmark from the v1 reduced-form analysis, and the gap is traced mechanically to two methodological additions here: UH-clean cost distributions and explicit pricing of tax distortion.

Third, a counterfactual 10\% SME price preference achieves comparable SME win-rate gains at essentially zero price cost. The welfare weight on SME surplus needed to rationalize the full set-aside over the price preference is 2.4--3.0, well above the 1.2--1.5 that Brazilian cash-transfer programs operate under. Under any plausible aggregator, the full set-aside is Pareto-inferior to a properly designed preference rule.

The policy implication for the implementation of Brazil's new procurement framework (Law 14{,}133/2021) is direct. The law contemplates both set-aside and price-preference instruments; my results suggest procurement authorities should lean toward the latter for non-specialized items. For pharma, where welfare costs are largest and the SME participation gain is smallest, the case for relaxation is strongest. For non-pharma, where the preference rule may be administratively more complex to implement, the choice is closer; but the point estimates still favor preference over set-aside by a wide margin.

The structural framework also suggests extensions worth pursuing. A first-price sealed-bid equilibrium model with preference rules would sharpen the V3 comparison. A dynamic model of SME capacity accumulation would price the longer-run benefits that the static analysis abstracts from. A multi-market extension would test whether the Pareto-superiority of the preference rule survives general-equilibrium feedback through input markets. All three are open.

The reduced-form fiscal-cost estimate in v1 left the structural interpretation open. The model and counterfactuals here close that gap: set-asides restructure bidding more than they restrict entry, they cost more in welfare once tax distortion is priced, and a simpler preference rule likely captures most of the protective value at a fraction of the cost.
